\chapter{Cartan Subalgebras and Roots}
\label{ch:roots}

Throughout this chapter \(\mathfrak g\) is a finite-dimensional
semisimple complex Lie algebra and \(B\) is its Killing form.
The words semisimple and nilpotent, when applied to an element, refer
to its adjoint endomorphism. Inner automorphisms are generated by
\(\exp(\operatorname{ad}x)\), \(x\in\mathfrak g\); the exponential here
is an ordinary convergent matrix exponential. No integration theorem
for abstract Lie algebras is needed.

\section{Tori and their centralizers}

\begin{proposition}[Maximal toral subalgebras]\label{res:D01}
A subalgebra consisting of commuting semisimple elements is called
toral. A maximal toral subalgebra \(\mathfrak h\) exists and satisfies
\[
 C_{\mathfrak g}(\mathfrak h)=\mathfrak h.
\]
In particular \(\mathfrak h\) consists of simultaneously diagonalizable
adjoint endomorphisms.
\end{proposition}
\begin{proof}
Commuting diagonalizable matrices are simultaneously diagonalizable:
an eigenspace of one is preserved by all the others, and induction on
dimension finishes the assertion. Consequently the sum of two commuting
toral subalgebras is toral. Maximal ones exist by finite dimensionality.
If \(\mathfrak g\ne0\), one is nonzero: otherwise every element would
be adjoint-nilpotent by the Jordan decomposition
of Theorem~\ref{res:B07}, and Engel would make \(\mathfrak g\) nilpotent.

Write \(\mathfrak c=C_{\mathfrak g}(\mathfrak h)\). Simultaneous
diagonalization gives
\(\mathfrak g=\bigoplus_\lambda\mathfrak g_\lambda\), where
\([h,x]=\lambda(h)x\) on \(\mathfrak g_\lambda\).
Invariance says
\[
 (\lambda(h)+\mu(h))B(x,y)=0
 \quad(x\in\mathfrak g_\lambda,\ y\in\mathfrak g_\mu).
\]
Thus \(B|_{\mathfrak c}\) is nondegenerate: a vector in its radical
is also orthogonal to every nonzero weight space and hence to all
of \(\mathfrak g\).

For \(x\in\mathfrak c\), both Jordan components \(x_s,x_n\) centralize
\(\mathfrak h\). This follows either from uniqueness of Jordan
decomposition or from the polynomial expressions for the two adjoint
components. Maximality implies \(x_s\in\mathfrak h\).
Therefore \(\operatorname{ad}_{\mathfrak c}x=
\operatorname{ad}_{\mathfrak c}x_n\) is nilpotent.
Engel implies that \(\mathfrak c\) is nilpotent, in particular solvable.
Lie's theorem triangularizes its action on \(\mathfrak g\).
The operators from \([\mathfrak c,\mathfrak c]\) then have zero
diagonal, so
\(B([\mathfrak c,\mathfrak c],\mathfrak c)=0\).
Nondegeneracy gives \([\mathfrak c,\mathfrak c]=0\).
Finally \(\operatorname{ad}x_n\) commutes with every
\(\operatorname{ad}y\), \(y\in\mathfrak c\). Their product is nilpotent,
so \(B(x_n,y)=0\) for every such \(y\). Hence \(x_n=0\).
Every element of \(\mathfrak c\) is semisimple and
\(\mathfrak c=\mathfrak h\).
\end{proof}

\begin{proposition}[Root decomposition and opposite pairings]\label{res:D02}
For a maximal toral \(\mathfrak h\) there is a decomposition
\[
 \mathfrak g=\mathfrak h\oplus
       \bigoplus_{\alpha\in R}\mathfrak g_\alpha,\qquad
 \mathfrak g_\alpha=\{x:[h,x]=\alpha(h)x\text{ for all }h\in\mathfrak h\},
\]
where \(R\subset\mathfrak h^*\setminus\{0\}\) is finite.
One has
\([\,\mathfrak g_\alpha,\mathfrak g_\beta\,]\subset
\mathfrak g_{\alpha+\beta}\), with the zero weight interpreted as
\(\mathfrak h\). Moreover \(B(\mathfrak g_\alpha,\mathfrak g_\beta)=0\)
unless \(\alpha+\beta=0\), and
\(\mathfrak g_\alpha\times\mathfrak g_{-\alpha}\to\mathbb C\)
is a perfect pairing.
\end{proposition}
\begin{proof}
The decomposition and its zero space follow from
Proposition~\ref{res:D01}. Jacobi gives
\([h,[x,y]]=(\alpha+\beta)(h)[x,y]\).
The orthogonality equation in that proof gives the asserted
orthogonality. If a vector in \(\mathfrak g_\alpha\) pairs to zero
with \(\mathfrak g_{-\alpha}\), it pairs to zero with every summand,
so is zero. Apply this on both sides of the finite-dimensional pairing.
The same argument establishes nondegeneracy on \(\mathfrak h\).
\end{proof}

\section{Root triples and root strings}

Use \(B|_{\mathfrak h}\) to define \(t_\alpha\in\mathfrak h\) by
\(B(t_\alpha,h)=\alpha(h)\), and put
\((\alpha,\beta)=B(t_\alpha,t_\beta)\). Invariance immediately gives
\[
 [e,f]=B(e,f)t_\alpha
 \quad(e\in\mathfrak g_\alpha,\ f\in\mathfrak g_{-\alpha}).
 \tag{4.1}\label{eq:opposite-bracket}
\]

\begin{theorem}[Root triples and reducedness]\label{res:D03}
For every root \(\alpha\), \((\alpha,\alpha)\ne0\).
There are \(e_\alpha\in\mathfrak g_\alpha\),
\(f_\alpha\in\mathfrak g_{-\alpha}\), and
\[
 h_\alpha=\frac{2t_\alpha}{(\alpha,\alpha)}
\]
with \([e_\alpha,f_\alpha]=h_\alpha\),
\([h_\alpha,e_\alpha]=2e_\alpha\), and
\([h_\alpha,f_\alpha]=-2f_\alpha\).
Every root space is one dimensional, and the only roots proportional
to \(\alpha\) are \(\alpha\) and \(-\alpha\).
\end{theorem}
\begin{proof}
Choose \(e,f\) with \(B(e,f)\ne0\), and let \(h=[e,f]\).
If \(\alpha(t_\alpha)=0\), the three-dimensional span of \(e,f,h\)
is solvable: \(h\) commutes with \(e,f\). Lie's theorem makes
\(\operatorname{ad}h=[\operatorname{ad}e,\operatorname{ad}f]\)
strictly upper triangular, hence nilpotent. Since \(h\in\mathfrak h\),
it is also semisimple. The zero center of \(\mathfrak g\) would
give \(h=0\), contradicting (4.1). Normalize \(B(e,f)\) to
\(2/(\alpha,\alpha)\); (4.1) gives the three relations.

First consider the root spaces whose roots are integral multiples
of \(\alpha\), together with \(\mathbb C h_\alpha\).
Their sum is a subalgebra stable under this \(\mathfrak{sl}_2\).
All its \(h_\alpha\)-weights are even, and its zero space is
one dimensional. Complete reducibility for \(\mathfrak{sl}_2\)
implies that it is one simple even-highest-weight module: every
such simple summand has a one-dimensional zero space.
It contains the adjoint three-dimensional module
\(\langle e_\alpha,h_\alpha,f_\alpha\rangle\), so equals that module.
Thus \(\dim\mathfrak g_\alpha=1\) and no integral multiple
\(m\alpha\), \(|m|\ge2\), is a root.

For an arbitrary proportional root \(c\alpha\), integrality of
\(\mathfrak{sl}_2\)-weights gives \(2c\in\mathbb Z\).
Interchanging the two roots gives \(2/c\in\mathbb Z\).
Consequently \(c\in\{\pm1,\pm2,\pm\tfrac12\}\).
The last four cases contradict the already established exclusion
of integral multiples, applied to one of the two roots.
\end{proof}

\begin{theorem}[Strings, integers, and reflections]\label{res:D04}
If \(\beta\ne\pm\alpha\), the roots in \(\beta+\mathbb Z\alpha\)
are exactly
\[
 \beta-p\alpha,\ldots,\beta+q\alpha,\qquad
 p,q\ge0,\qquad p-q=\beta(h_\alpha).
\]
All successive nontrivial brackets along this string are nonzero.
In particular
\[
 \beta(h_\alpha)=\frac{2(\beta,\alpha)}{(\alpha,\alpha)}\in\mathbb Z,
 \qquad
 s_\alpha(\beta)=\beta-\beta(h_\alpha)\alpha\in R.
\]
The reflection is implemented on \(\mathfrak g\) by
\[
 n_\alpha=
 \exp(\operatorname{ad}e_\alpha)
 \exp(-\operatorname{ad}f_\alpha)
 \exp(\operatorname{ad}e_\alpha).
 \tag{4.2}\label{eq:root-reflection-lift}
\]
\end{theorem}
\begin{proof}
The finite sum
\(V=\bigoplus_{k\in\mathbb Z}\mathfrak g_{\beta+k\alpha}\)
is a module for the root triple. Each weight space for
\(h_\alpha\) is one dimensional, and all its weights differ by
even integers. Two nonzero simple summands would share the
weight \(0\) if their highest weights are even, or the weights
\(\pm1\) if odd. Hence \(V\) is simple.
Its \(\mathfrak{sl}_2\) weights form one interval, with endpoints
\(\beta(h_\alpha)-2p\) and \(\beta(h_\alpha)+2q\), opposite to each
other. This proves the equation and the nonzero bracket assertion.
Symmetry of that interval sends the root \(\beta\) to
\(\beta-\beta(h_\alpha)\alpha\).

On each finite simple \(\mathfrak{sl}_2\) module the three
exponentials in (4.2) interchange opposite weights. On the triple
itself direct multiplication gives \(h_\alpha\mapsto-h_\alpha\).
For \(h\in\ker\alpha\), all three factors fix \(h\), so
\(n_\alpha(h)=h-\alpha(h)h_\alpha\).
It follows by applying brackets that \(n_\alpha\) maps the entire
\(\beta\)-space to the \(s_\alpha\beta\)-space.
The exponentials are finite on these modules, and an exponential
of a derivation preserves brackets; thus they are automorphisms.
\end{proof}

\begin{proposition}[The real root space]\label{res:D05}
The roots span \(\mathfrak h^*\). Let
\(\mathfrak h_{\mathbb R}=\operatorname{span}_{\mathbb R}
\{h_\alpha:\alpha\in R\}\). Then
\[
 \mathfrak h=\mathfrak h_{\mathbb R}\oplus
 i\mathfrak h_{\mathbb R},
\]
and \(B\) is real positive definite on \(\mathfrak h_{\mathbb R}\).
The roots, viewed in its real dual \(E\), form a finite reduced
crystallographic root system. Their inner products for the
Killing-induced form are rational, and their rational span is
a rational form of \(E\).
\end{proposition}
\begin{proof}
If all roots vanish on \(h\), then \(\operatorname{ad}h=0\);
therefore \(h=0\). Nondegeneracy identifies the spanning property
with the fact that the \(h_\alpha\) span \(\mathfrak h\) over
\(\mathbb C\).
For \(h\) in their real span all eigenvalues of \(\operatorname{ad}h\)
are real, by Theorem~\ref{res:D04}. If also \(h\in
i\mathfrak h_{\mathbb R}\), they are purely imaginary and hence
zero, giving \(h=0\). Moreover
\[
 B(h,h)=\sum_{\beta\in R}\beta(h)^2>0\quad(h\ne0),
\]
and polarization makes \(B\) real valued on this real space.
Dimension now gives the displayed complexification.
The dual Euclidean form is the restriction of the form already
used for roots. Reducedness, the integral coroot pairings,
reflection invariance, and spanning have all been proved.
These are precisely the root system axioms.

The rationality assertion follows without choosing a Dynkin type.
By the one-dimensional root spaces and integral coroot pairings,
\[
 B(h_\alpha,h_\beta)
     =\sum_{\gamma\in R}\gamma(h_\alpha)\gamma(h_\beta)
       \in\mathbb Z.
\]
Since \(B(h_\alpha,h_\alpha)=4/(\alpha,\alpha)\ne0\),
the square length \((\alpha,\alpha)\) is rational. The identity
\((\alpha,\beta)=\beta(h_\alpha)(\alpha,\alpha)/2\)
then gives all rational inner products. Choose any real basis
from the roots. Its nonsingular Gram matrix has rational entries,
as do the pairings of any other root with that basis; solving
the Gram system expresses that root with rational coefficients.
Thus the rational span has dimension \(\dim_{\mathbb R}E\),
and its scalar extension to \(\mathbb R\) is \(E\).
\end{proof}

\section{Regular elements and conjugacy}

The following argument separates the ordinary kernel dimension
from the multiplicity of zero in a characteristic polynomial.
For \(x\in\mathfrak g\), let \(n(x)\) be the latter multiplicity,
and set \(r=\min_x n(x)\).

\begin{theorem}[Regular semisimple density]\label{res:D06}
The set of elements with \(n(x)=r\) is a connected dense open
subset of \(\mathfrak g\); every one of these elements is semisimple.
For every maximal toral \(\mathfrak h\),
\(\dim\mathfrak h=r\), and
\[
 \mathfrak h_{\rm reg}
 =\{h\in\mathfrak h:\alpha(h)\ne0\text{ for all }\alpha\in R\}.
\]
A regular semisimple element has a unique maximal toral
subalgebra containing it, namely its centralizer.
For every \(x\in\mathfrak g\),
\(\dim C_{\mathfrak g}(x)\ge r\), with equality on this open set.
\end{theorem}
\begin{proof}
For \(\mathfrak g=0\) every assertion is immediate, with
\(r=0\); hence assume \(\mathfrak g\ne0\).
Write
\(\det(t-\operatorname{ad}x)=t^r
\bigl(t^{m}+a_{m-1}(x)t^{m-1}+\cdots+a_0(x)\bigr)\),
where \(m=\dim\mathfrak g-r\).
By the definition of \(r\), \(a_0\) is a nonzero polynomial,
and its nonvanishing is exactly \(n(x)=r\).
The complement of a nonzero complex polynomial's zero set is
open and dense: vanishing on a ball forces all Taylor
coefficients to vanish. It is path connected because the complex
line between two of its points meets the zero set at only
finitely many parameters, and a path in that line can avoid them.

Fix a maximal toral \(\mathfrak h\), and put
\(V=\mathfrak h_{\rm reg}\) as defined in the statement.
For \(h\in V\),
\(\mathfrak g=[\mathfrak g,h]\oplus\mathfrak h\).
Choose a linear complement \(M\) to \(\mathfrak h\).
The analytic map
\[
 M\times\mathfrak h\longrightarrow\mathfrak g,\qquad
 (y,z)\longmapsto\exp(\operatorname{ad}y)z
\]
has differential \((y,z)\mapsto[y,h]+z\) at \((0,h)\),
an isomorphism. The inverse function theorem makes the inner
saturation of \(V\) open. It therefore meets the dense set
\(n(x)=r\). Nullity is invariant under automorphisms, and equals
\(\dim\mathfrak h\) on \(V\). Thus \(\dim\mathfrak h=r\).

Let first \(x\) be semisimple with \(n(x)=r\).
Its centralizer \(\mathfrak c\) has dimension \(r\), and
\(B|_{\mathfrak c}\) is nondegenerate by the zero weight argument.
For \(y\in\mathfrak c\), \(\operatorname{ad}(x+ty)\)
is invertible on the sum of nonzero \(x\)-eigenspaces for small
\(t\). On \(\mathfrak c\) it equals \(t\operatorname{ad}y\).
Since its zero multiplicity cannot be less than \(r=\dim\mathfrak c\),
\(\operatorname{ad}_{\mathfrak c}y\) is nilpotent.
Engel and Lie's theorem, followed by nondegeneracy of
\(B|_{\mathfrak c}\), show that \(\mathfrak c\) is abelian,
exactly as in Proposition~\ref{res:D01}.
For the same reason every nilpotent component of an element of
\(\mathfrak c\) is zero. Thus \(\mathfrak c\) is toral.

For arbitrary \(x\) of minimal nullity, the Jordan components
commute, and \(\operatorname{ad}x\) and \(\operatorname{ad}x_s\)
have the same generalized zero space. Thus \(n(x_s)=r\).
The preceding paragraph makes \(C_{\mathfrak g}(x_s)\) toral.
It contains \(x\), so \(x\) is semisimple and \(x_n=0\).
Any toral subalgebra containing \(x\) is contained in this
centralizer; maximality then gives uniqueness.

On the dense open set just described,
\(\operatorname{ad}x\) has rank \(\dim\mathfrak g-r\).
Every larger matrix minor therefore vanishes on a dense set
and hence identically. This proves the claimed lower bound
for the kernel dimension at all \(x\), not just semisimple ones.
\end{proof}

\begin{theorem}[Cartan subalgebras and their conjugacy]\label{res:D07}
A Cartan subalgebra, defined as a nilpotent self-normalizing
subalgebra, is the same as a maximal toral subalgebra of
\(\mathfrak g\). All of them are conjugate by inner automorphisms.
\end{theorem}
\begin{proof}
A maximal toral subalgebra is abelian and self-normalizing:
in the root decomposition, a nonzero root component cannot
normalize its zero space. It is therefore Cartan.

Conversely let \(\mathfrak a\) be nilpotent and
\(N_{\mathfrak g}(\mathfrak a)=\mathfrak a\).
We shall use the following elementary consequence of Engel:
if \(\mathfrak a\) acts nilpotently on a space \(V\supset\mathfrak a\)
and preserves \(\mathfrak a\), then \(V/\mathfrak a\ne0\)
has a nonzero common zero vector. For the adjoint action such
a vector lifts to a normalizing vector. Thus self-normalization
forces \(V=\mathfrak a\).

For \(x\in\mathfrak a\), both \(\operatorname{ad}x_s\) and
\(\operatorname{ad}x_n\), being polynomials in
\(\operatorname{ad}x\), preserve \(\mathfrak a\).
Hence \(x_s,x_n\in\mathfrak a\).
On \(\mathfrak a\) the action of \(x\) is nilpotent,
so the semisimple restriction of \(\operatorname{ad}x_s\)
is zero. The semisimple elements of \(\mathfrak a\) therefore
form a central torus \(\mathfrak t\).
On \(V=C_{\mathfrak g}(\mathfrak t)\) every \(x\in\mathfrak a\)
acts as the nilpotent operator \(\operatorname{ad}x_n\).
The preceding consequence of Engel gives \(V=\mathfrak a\).
In particular \(B|_{\mathfrak a}\) is nondegenerate.
Lie's theorem triangularizes \(\mathfrak a\) on \(\mathfrak g\);
the nilpotent components \(x_n\) have zero diagonal, hence
\(B(x_n,\mathfrak a)=0\). Thus all \(x_n=0\),
and \(\mathfrak a=\mathfrak t\) is maximal toral.

For conjugacy, every regular semisimple \(x\) has the unique
Cartan \(C(x)\). Declare two such elements equivalent if their
Cartans are inner-conjugate. For any fixed Cartan
\(\mathfrak h\), the class in question is precisely its inner
saturation of \(\mathfrak h_{\rm reg}\).
The differential calculation in Theorem~\ref{res:D06}
makes each such class open. The regular semisimple set is
connected, so its partition into open classes has a single class.
Every Cartan contains a regular element; therefore all Cartans
are conjugate.
\end{proof}

\section{Matrix models}

\begin{example}[The classical root systems]\label{res:D08}
In the following table the \(\varepsilon_i\) are coordinate
functionals on diagonal Cartans:
\[
\begin{array}{c|c|c}
\mathfrak g&\mathfrak h&R\\ \hline
\mathfrak{sl}_{n+1}&
 \operatorname{diag}(t_1,\ldots,t_{n+1}),\ \sum t_i=0&
 \varepsilon_i-\varepsilon_j\ (i\ne j)\\
\mathfrak{so}_{2n+1}&
 \operatorname{diag}(t_1,\ldots,t_n,0,-t_1,\ldots,-t_n)&
 \pm\varepsilon_i,\ \pm\varepsilon_i\pm\varepsilon_j\\
\mathfrak{sp}_{2n}&
 \operatorname{diag}(t_1,\ldots,t_n,-t_1,\ldots,-t_n)&
 \pm2\varepsilon_i,\ \pm\varepsilon_i\pm\varepsilon_j\\
\mathfrak{so}_{2n}&
 \operatorname{diag}(t_1,\ldots,t_n,-t_1,\ldots,-t_n)&
 \pm\varepsilon_i\pm\varepsilon_j .
\end{array}
\]
The orthogonal forms are chosen in hyperbolic coordinates.
Coroots are \(2\alpha/(\alpha,\alpha)\) under the Euclidean
identification. The low rank coincidences are
\(B_1=C_1=A_1\), \(B_2=C_2\), \(D_2=A_1\sqcup A_1\),
and \(D_3=A_3\); \(D_1\) is not a semisimple case.
\end{example}
\begin{proof}
For a diagonal matrix \(t\),
\([t,E_{ij}]=(t_i-t_j)E_{ij}\).
In the symplectic case the equation \(X^TJ+JX=0\) gives
\[
 X=\begin{pmatrix}A&B\\C&-A^T\end{pmatrix},
 \qquad B=B^T,\quad C=C^T.
\]
The off-diagonal entries of \(A\) have weights
\(\varepsilon_i-\varepsilon_j\); the off-diagonal symmetric
entries of \(B,C\) have weights
\(\pm(\varepsilon_i+\varepsilon_j)\), and their diagonal
entries have weights \(\pm2\varepsilon_i\).
For the even orthogonal form
\(\left(\begin{smallmatrix}0&I\\I&0\end{smallmatrix}\right)\),
the same block equation makes \(B,C\) skew symmetric and
therefore removes the diagonal weights.
In odd dimension the extra row and column give
\(\pm\varepsilon_i\). These calculations account for every
matrix entry allowed by the defining equations, so the zero
spaces are exactly the displayed Cartans.
Pair opposite matrix entries to obtain the corresponding
root triples, which gives the stated coroot normalization.
Finally the listed small root sets are carried into each
other by permuting or rescaling the Euclidean coordinates;
their Lie algebra isomorphisms also follow from the
reconstruction theorem proved in Chapter~\ref{ch:classification}.
\end{proof}

The algebraic root calculations follow
\cite[Sections 8 and 15--16]{humphreys} and
\cite[Sections 19--20]{etingof}. The use of the characteristic
polynomial and its connected nonvanishing locus in the conjugacy
argument explains why the assertion about a dense regular set
concerns regular \emph{semisimple} elements, not every element
whose ordinary centralizer has minimal dimension.
